% Generated draft: requires mathematical and source review.
\documentclass[11pt]{article}
\usepackage[margin=25mm]{geometry}
\usepackage{fontspec}
\setmainfont{Latin Modern Roman}
\usepackage{amsmath,amssymb,mathtools,unicode-math}
\setmathfont{Latin Modern Math}
\usepackage{enumitem,needspace,xurl,hyperref}
\setlength{\parindent}{0pt}
\setlength{\parskip}{5pt}
\setlength{\emergencystretch}{3em}
\newcommand{\N}{\mathbb{N}}
\newcommand{\Z}{\mathbb{Z}}
\newcommand{\Q}{\mathbb{Q}}
\newcommand{\R}{\mathbb{R}}
\newcommand{\C}{\mathbb{C}}
\newcommand{\F}{\mathbb{F}}
\newcommand{\A}{\mathbb{A}}
\newcommand{\PP}{\mathbb{P}}
\newcommand{\E}{\mathbb{E}}
\newcommand{\eps}{\varepsilon}
\newcommand{\dd}{\,\mathrm{d}}
\begin{document}
\section*{UnsolvedMath ID: 30004111}
\textbf{Source problem code:} catalogue entry
\subsection{Definitions and mathematical statement}
For a finite group \(G\) and a prime \(\ell\), write \(\operatorname{Irr}(G)\) for its ordinary irreducible characters and \(\operatorname{IBr}_{\ell}(G)\) for its irreducible Brauer characters in characteristic \(\ell\), over a splitting field. If \(\chi^{\circ}\) denotes the restriction of \(\chi\in\operatorname{Irr}(G)\) to the elements whose orders are prime to \(\ell\), the decomposition numbers \(d_{\chi\varphi}\in\mathbb Z_{\geq 0}\) are defined by
\[
\chi^{\circ}=\sum_{\varphi\in\operatorname{IBr}_{\ell}(G)}d_{\chi\varphi}\varphi.
\]
The corresponding matrix is the \(\ell\)-decomposition matrix.

\textbf{Supplied conjectural target (Geck's Unitriangularity Conjecture).} For every finite reductive group \(G\) and every prime \(\ell\) that is very good in the intended reductive setting, can the unipotent decomposition matrix be ordered by unipotent classes so that it is block lower triangular, with identity diagonal blocks indexed by the relevant families of unipotent characters?

The matrix condition can be made explicit once the intended rows, columns, families, and class order have been specified. Let \(U\subseteq\operatorname{Irr}(G)\) be the intended unipotent character row set and let \(B\subseteq\operatorname{IBr}_{\ell}(G)\) be the intended column set. In a square interpretation, the requested condition is the existence of ordered partitions
\[
U=\bigsqcup_{i=1}^{r}U_i,\qquad
B=\bigsqcup_{i=1}^{r}B_i,\qquad
|U_i|=|B_i|=n_i,
\]
with enumerations \(U_i=\{\chi_{i,1},\ldots,\chi_{i,n_i}\}\) and \(B_i=\{\varphi_{i,1},\ldots,\varphi_{i,n_i}\}\), such that
\[
d_{\chi_{i,a},\varphi_{j,b}}=0\quad(i<j),
\qquad
d_{\chi_{i,a},\varphi_{i,b}}=\delta_{ab}.
\]
Here \(\delta_{ab}\) is the Kronecker delta. The row blocks must be the intended unipotent-character families, and their ordering must have the intended compatibility with unipotent classes.

\textbf{Unresolved qualifications.} The supplied record does not specify the reductive realization of \(G\), the root datum relative to which \(\ell\) is very good, whether nondefining characteristic is required, the precise unipotent row and Brauer column sets, the meaning of the relevant families, or the association and ordering of unipotent classes. It also does not say whether the matrix is square, is a selected square submatrix, or has additional rows or columns. Thus the displayed square formulation explains the matrix condition conditionally; it does not select a submatrix or resolve the missing conventions. The source's background discusses remaining bad-characteristic cases, whereas the stated target assumes a very good prime; the relationship between these scopes is unresolved.
\subsection{Short English statement}
The supplied conjecture asks whether, for every finite reductive group and every very good prime \(\ell\), the unipotent decomposition matrix can be arranged by unipotent classes into lower triangular blocks, with identity blocks corresponding to unipotent-character families. The precise matrix, families, and class order require clarification. The supplied partially solved label is dated curation, not an independently verified conclusion.
\subsection{Sources}
\begin{itemize}
\item Local UnsolvedMath registry, supplied complete record: permanent ID \(30004111\), exact source problem code \texttt{catalogue entry}, title \emph{Geck's Unitriangularity Conjecture}, collection \emph{Oberwolfach Reports Open Problems}. The conjecture is attributed by name to Geck; no proposer is separately identified. The supplied proposed year is 2019.
\item Supplied report citation: \emph{Representations of Finite Groups (2020)}, Oberwolfach Reports, pp.~841--895, DOI \texttt{10.4171/owr/2019/14}. \url{https://doi.org/10.4171/owr/2019/14}. No exact problem page or passage is supplied.
\item Supplied literature reference, labelled ``Geck unitriangularity theorem'': \url{https://doi.org/10.4007/annals.2020.192.2.7}.
\item Supplied literature reference, labelled ``Extensions of unitriangularity'': \url{https://arxiv.org/abs/2402.17616}.
\item Supplied literature reference, labelled ``Oberwolfach representation theory report'': \url{https://doi.org/10.4171/owr/2019/14}.
\item Literature assessment dated \texttt{2026-08-22}: the supplied curation labels the problem partially solved and states that Geck's unitriangularity conjecture is proved in broad good-characteristic settings, while the general extracted conjecture, including remaining bad-characteristic cases, is not fully resolved. This is supplied dated curation derived from the OWR record, not an independently verified assessment. Its scope relative to the very-good-prime hypothesis requires review.
\item The supplied statement status is \texttt{corrected\_verified}, with a recorded manual semantic, typographic, and LaTeX review claiming that the clean statement is self-contained and preserves the target. This is source metadata; the present draft does not independently verify that claim or the references.
\end{itemize}
\end{document}
